% !TeX encoding = UTF-8
% !TeX spellcheck = de_DE
% Datei: 02_theoreme.tex
% Diese Datei enthält NUR den Kapitelinhalt und wird via \input in main.tex eingebunden.

\chapter{Theoretische Ableitungen und Theoreme}
\label{chap:theoreme}

\section{Einleitung}
Aus dem axiomatischen System $\mathcal{S}$ (Definiert in Abschnitt~\ref{sec:synthese}) lassen sich zwingende logische und strukturelle Konsequenzen für jede Gesellschaft $G$ ableiten, die $\mathcal{S}$ instanziiert. Diese Konsequenzen werden hier als Theoreme formalisiert. Sie beschreiben die notwendige Beschaffenheit einer $\mathcal{S}$-Gesellschaft jenseits der Grundpostulate und bilden die Brücke zur angewandten Institutionsgestaltung und Stabilitätsanalyse.

\section{Grundlegende Definitionen}
Zur präzisen Formulierung der Theoreme werden folgende Begriffe definiert:

\begin{definition}[Demokratische Infrastruktur]
    Die demokratische Infrastruktur $\mathit{DI}(G)$ einer Gesellschaft $G$ ist die Gesamtheit aller Einrichtungen und Verfahren, die der Informationsbeschaffung, Deliberation und Entscheidungsfindung gemäß Axiom~\ref{ax:4} dienen. Formal ist $\mathit{DI}(G)$ ein Tupel aus Bildungs-, Medien-, Beratungs- und Entscheidungssystemen.
\end{definition}

\begin{definition}[Ökonomisches Gleichgewichtsintervall]
    Sei $\varepsilon$ die Reichtumsgrenze aus Axiom~\ref{ax:3}. Das ökonomische Gleichgewichtsintervall $I_{\mathit{eq}} \subseteq \mathbb{R}^+$ ist dasjenige Intervall um $\varepsilon$, innerhalb dessen die Stabilität von $\mathcal{S}$ gewahrt bleibt. Seine Existenz und Eigenschaften sind Gegenstand der Variationsrechnung (Kapitel~\ref{chap:variation}).
\end{definition}

\section{Die Haupttheoreme der $\mathcal{S}$-Gesellschaft}

\subsection{Theorem 1: Die Auflösung des Repräsentationsparadigmas}

\begin{theorem}[Struktur der Entscheidungsfindung]
    In einer $\mathcal{S}$-Gesellschaft $G$ existiert keine stabile politische Repräsentation im klassischen Sinne. Die Entscheidungsfindung folgt notwendig einem von zwei Modi:
    \begin{enumerate}
        \item \textbf{Direkte Entscheidung:} Für Angelegenheiten von überschaubarer Komplexität und Reichweite.
        \item \textbf{Strikt delegierte Entscheidung:} Für komplexe oder großskalige Angelegenheiten, wobei die Delegierten (\glqq Beauftragte\grqq) ein imperatives Mandat haben und jederzeit ohne Verfahrenshürden abberufbar sind.
    \end{enumerate}
\end{theorem}

\begin{proof}
    Direkt aus Axiom~\ref{ax:2} ($P(G)=\varnothing$) folgt, dass keine dauerhaften Repräsentationskörper mit eigenem Entscheidungsspielraum (freiem Mandat) existieren können. Die einzig verbleibenden, mit Axiom~\ref{ax:1} und \ref{ax:2} kompatiblen Organisationsformen sind die genannten: unmittelbare Partizipation oder extrem kontrollierte, reversible Delegation.
\end{proof}

\subsection{Theorem 2: Die Notwendigkeit der Bildungsgesellschaft}

\begin{theorem}[Bildung als zentrale Ressource]
    Sei $B(G) \in \mathbb{R}^+$ das Maß für die in das Bildungssystem von $G$ investierten Ressourcen (pro Kopf). Sei $\sigma(G)$ das Maß für die Informationsasymmetrie in $G$. Dann gilt für eine stabile $\mathcal{S}$-Gesellschaft:
    \[
    \boxed{\exists\, B_{\min} > 0: \quad B(G) < B_{\min} \quad \Rightarrow \quad \sigma(G) \text{ wächst, was} \quad \mathcal{D}(G) \to 0 \quad \text{bewirkt}.}
    \]
    Mit anderen Worten: Ein Unterschreiten einer kritischen Bildungsinvestition destabilisiert das System fundamental.
\end{theorem}

\begin{proof}
    Axiom~\ref{ax:4} fordert, dass Entscheidungen auf wissenschaftlichen Fakten basieren und von der Gemeinschaft verstanden und genehmigt werden. Eine ungenügende Bildung $B(G) < B_{\min}$ führt zu:
    \begin{enumerate}
        \item Einem Anstieg von $\sigma(G)$ (Wissenskluft).
        \item Der Unfähigkeit der Gemeinschaft, $E$ (den evidenzbasierten Maßnahmenkatalog) zu bewerten.
        \item Dadurch wird der Prozess $E \to Z(G)$ hinfällig oder degeneriert zu einem bloßen Akklamationsritual. Dies verletzt den Geist von Axiom~\ref{ax:4} und untergräbt letztlich die rationale Grundlage von $\mathcal{D}(G)$.
    \end{enumerate}
    Aus der Stabilitätsanalyse (Kapitel~\ref{chap:variation}, Modul Informationsasymmetrie) folgt die Existenz der unteren Schranke $B_{\min}$.
\end{proof}

\subsection{Theorem 3: Die strukturelle Verhinderung von Kapitaloligarchien}

\begin{theorem}[Zerfall der Kapitalmacht]
    In einer $\mathcal{S}$-Gesellschaft kann sich keine dauerhafte, machtausübende Elite auf der Basis von ökonomischem Kapital bilden. Die Konversionsrate $c$ von ökonomischer Potenz $V_i$ in politischen Einfluss ist strukturell auf Null gesetzt.
    \[
    \boxed{\forall i \in I: V_i \leq \varepsilon \quad \Rightarrow \quad c = 0.}
    \]
\end{theorem}

\begin{proof}
    Die absolute und durchsetzbare Obergrenze $\varepsilon$ (Axiom~\ref{ax:3}) verhindert die Akkumulation von massivem, anderweitig überwältigendem Privatvermögen. Da politischer Einfluss durch Medienkontrolle, Lobbyismus oder Bestechung typischerweise mit Ressourceneinsatz verbunden ist, wird diese Möglichkeit durch die Begrenzung von $V_i$ ausgeschlossen. Die Verfassung (Axiom~\ref{ax:5}) garantiert die Dauerhaftigkeit dieser Schranke.
\end{proof}

\subsection{Theorem 4: Wissenschaft als öffentliche und kontrollierte Infrastruktur}

\begin{theorem}[Organisation der Wissensproduktion]
    Die wissenschaftlichen Einrichtungen in einer $\mathcal{S}$-Gesellschaft $G$ können nicht privatwirtschaftlich oder autokratisch organisiert sein. Sie müssen als transparente, öffentlich kontrollierte und pluralistisch besetzte Infrastruktur $\mathit{FI} \subset \mathit{DI}(G)$ (Forschungsinfrastruktur) ausgestaltet sein.
\end{theorem}

\begin{proof}
    Folgt aus Axiom~\ref{ax:4} und Theorem 3.
    \begin{enumerate}
        \item \textbf{Öffentliche Kontrolle:} Private Kontrolle über die Wissensproduktion (z.B. durch Konzerne) würde eine indirekte Macht über den Input $E$ des Entscheidungsprozesses bedeuten, was Theorem 3 widerspricht und $\mathcal{D}(G)$ verringert.
        \item \textbf{Transparenz und Pluralismus:} Da Wissenschaft fehlbar ist (explizit in Axiom~\ref{ax:4} anerkannt), muss die Erstellung von $E$ vor Gruppendenken und Paradigmenblindheit geschützt werden. Dies erfordert institutionell gesicherten methodischen und theoretischen Pluralismus sowie vollständige Transparenz der Daten und Modelle.
    \end{enumerate}
\end{proof}

\subsection{Theorem 5: Beschränkte Souveränität}

\begin{theorem}[Grenzen der demokratischen Entscheidung]
    In einer $\mathcal{S}$-Gesellschaft $G$ ist die Souveränität der demokratischen Entscheidung $Z(G)$ absolut, aber \emph{inhaltlich beschränkt} durch die Verfassung $C$.
    \[
    \boxed{\forall \text{ Entscheidungen } d: \quad Z(G) \text{ kann } d \text{ nur dann legitimieren, wenn } d \text{ mit } C \text{ vereinbar ist}.}
    \]
    Insbesondere kann $Z(G)$ nicht legitim über die Abschaffung oder Aufweichung von Axiom~\ref{ax:1}, \ref{ax:2}, \ref{ax:3} oder der demokratischen Verfahren von \ref{ax:4} entscheiden.
\end{theorem}

\begin{proof}
    Direkte Folge aus Axiom~\ref{ax:5} (Unveränderlichkeit der Verfassung $C$). Die Verfassung operationalisiert die Axiome 1-4. Jede Entscheidung $d$, die diesen widerspricht, wäre ein impliziter Verfassungsbruch, der im System $\mathcal{S}$ keine legale Form haben kann. Die demokratische Souveränität wirkt somit ausschließlich \emph{innerhalb} des durch $C$ vorgegebenen Rahmens.
\end{proof}

\section{Synthese: Die $\mathcal{S}$-Gesellschaft als kohärentes System}

Die fünf Theoreme zeigen, dass aus den knappen Axiomen ein dichtes Netz von strukturellen Vorgaben erwächst:
\begin{itemize}
    \item Theorem 1 und 3 zerstören die klassischen Säulen der Macht: politische Repräsentation und ökonomische Oligarchie.
    \item Theorem 2 und 4 errichten die neuen Säulen der $\mathcal{S}$-Gesellschaft: ein allumfassendes Bildungssystem und eine demokratisch eingebettete Wissenschaft.
    \item Theorem 5 sichert das gesamte Gefüge vor seiner eigenen Auflösung durch demokratischen Beschluss.
\end{itemize}

Damit ist das Bild einer $\mathcal{S}$-Gesellschaft bereits deutlich konturierter: Es handelt sich um eine hochgebildete, wissenschaftsbasierte, direkt-deliberative Gemeinschaft mit strikter ökonomischer Gleichheitsgrenze, deren Grundregeln für alle Ewigkeit festgeschrieben sind.

Die folgenden Kapitel werden untersuchen, ob und unter welchen Bedingungen ein solches System stabil existieren kann (Variationsrechnung) und mit welchen praktischen und philosophischen Herausforderungen es konfrontiert ist (Diskussion).
